\documentclass[12pt,a4paper]{article}
\usepackage[margin=25mm, headheight=15pt, headsep=10pt]{geometry}
\usepackage{fontspec}
\usepackage[table]{xcolor} % Note: table option is required for rowcolors
\usepackage{titlesec}
\usepackage{tocloft}
\usepackage{fancyhdr}
\usepackage{booktabs}
\usepackage{tcolorbox}
\tcbuselibrary{skins, breakable}
\usepackage[colorlinks=true, linkcolor=emerald, urlcolor=emerald, bookmarks=true]{hyperref}

\definecolor{emerald}{RGB}{0,102,51}
\definecolor{darkred}{RGB}{139,0,0}
\definecolor{lightgray}{RGB}{245,245,245}

\usepackage[nil]{babel}
\babelprovide[import, main]{bengali}
\babelprovide[import]{arabic}
\babelfont{rm}[Renderer=HarfBuzz,Scale=1.0]{Noto Serif Bengali}
\babelfont{sf}[Renderer=HarfBuzz]{Noto Sans Bengali}
\babelfont{tt}[Renderer=HarfBuzz]{Noto Sans Bengali}
\babelfont[arabic]{rm}[Renderer=HarfBuzz,Scale=1.1]{Noto Naskh Arabic}

% Section styling
\titleformat{\section}{\Large\bfseries\color{emerald}}{\thesection.}{0.5em}{}
\titleformat{\subsection}{\large\bfseries\color{emerald!85!black}}{\thesubsection.}{0.5em}{}
\titleformat{\subsubsection}{\normalsize\bfseries\color{emerald!70!black}}{\thesubsubsection.}{0.5em}{}
\titlespacing*{\section}{0pt}{18pt}{8pt}

% Fancy Header/Footer setup
\pagestyle{fancy}
\fancyhf{} % clear all header and footer fields
\fancyhead[L]{\color{emerald}\small\textbf{\nouppercase{\leftmark}}}
\fancyfoot[L]{\scriptsize\color{black!50}{\lat{AI}}-সংকলিত, আলেম-যাচাইকৃত নয়}
\fancyfoot[C]{\color{emerald!70!black}\thepage}
\renewcommand{\headrulewidth}{0.4pt}
\renewcommand{\headrule}{\hbox to\headwidth{\color{emerald}\leaders\hrule height \headrulewidth\hfill}}
\renewcommand{\footrulewidth}{0pt}

% Custom boxes
% 1. Ayat/Hadith Box
\newtcolorbox{ayatbox}[1][]{
    enhanced, breakable, 
    colback=emerald!5, 
    colframe=emerald, 
    boxrule=0.5pt, 
    arc=3pt, 
    left=10pt, right=10pt, top=10pt, bottom=10pt,
    #1
}

% 2. Quote/Translation Box
\newtcolorbox{quotebox}[1][]{
    enhanced, breakable,
    frame hidden,
    colback=lightgray,
    borderline west={3pt}{0pt}{emerald},
    left=15pt, right=10pt, top=10pt, bottom=10pt,
    sharp corners,
    #1
}

% 3. AI-generated notice box
\newtcolorbox{warnbox}[1][]{
    enhanced, breakable,
    colback=red!4,
    colframe=darkred,
    boxrule=0.8pt,
    arc=3pt,
    left=10pt, right=10pt, top=10pt, bottom=10pt,
    #1
}

% Commands
\newcommand{\arab}[1]{\foreignlanguage{arabic}{#1}}
\newcommand{\kib}[1]{\textcolor{emerald}{\textbf{#1}}}

% Latin (English) fallback: Noto Serif Bengali has NO Latin glyphs
\newfontfamily\latfont[Renderer=HarfBuzz]{Noto Serif}
\newcommand{\lat}[1]{{\latfont #1}}

\setlength{\parskip}{5pt}
\setlength{\parindent}{0pt}

% Fix for missing bullet in Noto Serif Bengali
\renewcommand{\labelitemi}{$\bullet$}



\begin{document}
\begin{center}{\Large\bfseries\color{emerald} \lat{15}-ব্যবহারিক-গাইড-খুশু-জনন}\end{center}
\vspace{1em}
\section*{১৫ — ব্যবহারিক গাইড: খুশু জনন ও নামাজে মনোযোগ}

\begin{warnbox}
 \textbf{গুরুত্বপূর্ণ নোটিশ:} এই কন্টেন্টটি \lat{AI} দিয়ে প্রস্তুত। কেবল সহীহ/হাসান হাদিস রাখা হয়েছে।
\end{warnbox}

\medskip\hrule\medskip

\subsection*{১. খুশু কী?}

\textbf{\lat{Khushu} (খুশু)} মানে নামাজে \textbf{মনোযোগ}, \textbf{শ্রদ্ধা}, \textbf{বিনয়} ও \textbf{আত্মসমর্পণ}। এটি \lat{heart} (অন্তর), \lat{mind} (মন) ও \lat{soul} (প্রাণ)-এর \lat{connection} (সংযোগ) আল্লাহর সাথে।

\textbf{\lat{“}কে জানে আমার উম্মাহর মধ্যে কতজন এমন আছে যারা শুধু মুখ দিয়ে নামাজ পড়ে — তাদের অন্তর কিছুই থাকে না।\lat{”}} (সুনানে তিরমিযী ২৯৯২ — সহীহ)

\medskip\hrule\medskip

\subsection*{২. নামাজের আগে প্রস্তুতি}

\begin{table}[h]\centering\small
\begin{tabular}{ll}
ধাপ & কীভাবে করবেন \\
\hline
\textbf{ওয়ু} & প্রতিটি অঙ্গ ধুলে মনে করুন: "আমি পবিত্র হচ্ছি — আল্লাহর সামনে দাঁড়ানোর জন্য" \\
\textbf{পোশাক} & পরিষ্কার পোশাক; ইত্র ব্যবহার করতে পারেন — এটি সম্মান \\
\textbf{মসজিদ} & ধীরে হাঁটুন; ইমামের পেছনে দাঁড়ান \\
\textbf{নিয়ত} & মনে করুন: "আমি \lat{X} রাকাত ফরজ/সুন্নাত নামাজ পড়ব — আল্লাহর মুখ অনুসন্ধান করে" \\
\hline
\end{tabular}
\end{table}

\medskip\hrule\medskip

\subsection*{৩. নামাজের মধ্যে খুশু}

\begin{table}[h]\centering\small
\begin{tabular}{ll}
বিষয় & কীভাবে প্রয়োগ করবেন \\
\hline
\textbf{কিরাআত} & ধীরে পড়ুন; প্রতিটি আয়াতের অর্থ মনে করুন \\
\textbf{চোখ} & কিবলার দিকে; সিজদার জায়গায় চোখ রাখুন \\
\textbf{হাত-পা} & শিথিল রাখুন — \lat{tension} (উত্তেজনা) নয় \\
\textbf{শ্বাস} & নাকে শ্বাস নিন; ধীরে বের করুন \\
\hline
\end{tabular}
\end{table}

\medskip\hrule\medskip

\subsection*{৪. খুশু বাড়ানোর ১০টি কৌশল}

\begin{table}[h]\centering\small
\begin{tabular}{lll}
\# & কৌশল & বিবরণ \\
\hline
১ & ইতমিনান & তাড়াহুড়ো করবেন না — প্রতিটি মুহূর্ত উপভোগ করুন \\
২ & অর্থ জানুন & যা পড়ছেন তার অর্থ জানুন — এটি \lat{fundamental} (মৌলিক) \\
৩ & সময়মতো & প্রতিটি নামাজ সময়মতো পড়ুন \\
৪ & একাই পড়ুন & মাঝে মাঝে একাই নামাজ পড়ুন — \lat{connection} (সংযোগ) বাড়ে \\
৫ & রাতের নামাজ & তাহাজ্জুদ — \lat{spiritual} \lat{peak} (আধ্যাত্মিক উচ্চতা) \\
৬ & দোয়া & নামাজে আল্লাহর সাথে \lat{conversation} (সংলাপ) করুন \\
৭ & ওয়াকফা & প্রতিটি মুহূর্তে ২-৩ সেকেন্ড \lat{pause} (বিরতি) দিন \\
৮ & আঙুল & শাহাদাতের সময় মাত্র আঙুল উঁচু করুন \\
৯ & দৃঢ় সংকল্প & প্রতিটি নামাজের আগে \lat{new} \lat{intention} (নতুন নিয়ত) করুন \\
১০ & ভয় ও আশা & জাহান্নামের ভয় ও জান্নাতের আশা — দুটোই রাখুন \\
\hline
\end{tabular}
\end{table}

\medskip\hrule\medskip

\subsection*{৫. সাধারণ ভুল ও সংশোধন}

\begin{table}[h]\centering\small
\begin{tabular}{ll}
ভুল & সংশোধন \\
\hline
তাড়াহুড়ো & ধীরে পড়ুন — প্রতিটি শব্দের ওজন অনুভব করুন \\
চোখ খোলা & চোখ বন্ধ নয়, কিন্তু নিচে — সিজদার জায়গায় \\
হাত নাড়া & হাত স্থির রাখুন \\
কথা বলা & নামাজে অনাবশ্যক কথা বলা যাবে না \\
\hline
\end{tabular}
\end{table}

\medskip\hrule\medskip

\subsection*{৬. সাপ্তাহিক অনুশীলন পরিকল্পনা}

\begin{table}[h]\centering\small
\begin{tabular}{ll}
সপ্তাহ & ফোকাস \\
\hline
১ম সপ্তাহ & ফাতিহার অর্থ মুখস্থ করুন \\
২য় সপ্তাহ & রুকু-সিজদার তাসবিহ ও দোয়া মুখস্থ করুন \\
৩য় সপ্তাহ & তাশাহহুদ ও দরুদ মুখস্থ করুন \\
৪র্থ সপ্তাহ & প্রতিটি নামাজে ধীরে পড়ার অনুশীলন \\
\hline
\end{tabular}
\end{table}

\medskip\hrule\medskip

\subsection*{৭. সংশ্লিষ্ট হাদিস}

\begin{table}[h]\centering\small
\begin{tabular}{ll}
হাদিস & গ্রন্থ \\
\hline
\textbf{\lat{“}বান্দা তার রবের সবচেয়ে কাছে থাকে সিজদার অবস্থায়।\lat{”}} & সহীহ মুসলিম ৪৮২ \\
\textbf{\lat{“}নামাজ পড়ো, যেমন আমাকে নামাজ পড়তে দেখেছ।\lat{”}} & সহীহ বুখারি ৬৩১ \\
\textbf{\lat{“}জামাতে ২৭ গুনাহ কমে যায়।\lat{”}} & সহীহ বুখারি ৬৪৫ \\
\hline
\end{tabular}
\end{table}

\medskip\hrule\medskip

\textit{শেষ আপডেট: আগস্ট ২০২৬}

\end{document}
